\subsection{Analysis plan, amendments and inference}
\label{sec:inference}

\paragraph{Pre-specified primary analysis.}
The ridge-readout analysis (Section~\ref{sec:ridge}) followed a written analysis plan
with a frozen configuration (schema version~2, configuration hash \texttt{206fa875cca5}),
dated 20~August~2026. The plan was not publicly registered; its dating rests on in-text
date stamps and the configuration hash, and the files were placed under version control
on 3~October~2026. We therefore describe its criteria as fixed before the analysis rather
than as preregistered. The plan asked whether the untrained-minus-backpropagation
difference at V1 reported in \citet{evalres} survives a fitted readout on the full Conv1
feature map. It fixed the layer, the voxel selection (split-half reliability above~0.2),
the grid ladder $g \in \{1, 2, 4, 8\}$, the cross-validation scheme and seed, the
regularisation grid and the aggregation (median over voxels, mean over subjects). The
unit of inference was the training seed ($n = 5$), tested with a one-sided exact
sign-flip test over all 32 sign assignments. The effect was to count as surviving only
if it held at $g \geq 4$. The minimal detectable effect $\Delta_{\min}$ was derived from
the pilot's seed-level standard deviation by simulating this test. The reference effect
was the value reported in version~1 of \citet{evalres}
($+\PreregReferenceEffect \pm \PreregReferenceSem$ at 224\,px); the current value is
$+\VoneGapRawTwoTwoFour$ [$\VoneGapRawTwoTwoFourLo$, $\VoneGapRawTwoTwoFourHi$], and the verdict does not depend on this choice because all
observed contrasts were negative (Section~\ref{sec:signflip}).

\paragraph{Amendments.}
Three amendments were adopted before the pilot on real data
% VERIFY: confirm in PHASE2_REPORT / PILOT_REPORT that D1-D3 preceded the real-data pilot
: a hard assertion against selecting the lowest regularisation value for reliable
voxels, selection of the regularisation by mean squared error with Pearson~$r$ reported,
and a widening of the regularisation grid. Three further decisions were taken after data
had been seen, and we list them as such. (i) The regularisation was kept at
minimum-MSE selection rather than a tolerance rule, after a tolerance measurement was
known; this adds no degree of freedom. (ii) The assertion's abort was replaced by a
post-hoc gate that classifies each lower-boundary selection as a plateau or a true
optimum below the grid; only the latter would invalidate a cell. This was decided after
a diagnosis showed that all observed boundary selections were plateaus, and it leaves all
scores unchanged. No cell was excluded in any analysis. (iii) Features for $g = 2$ were
extracted on CPU, whereas all other features came from a GPU extraction. The
pre-specified reproduction gate (absolute tolerance $10^{-6}$) failed because of
floating-point differences; we replaced it by a downstream-equivalence gate, which
passed (maximum absolute difference in median $r$ of \GtwoGateMaxDelta{} at $g = 4$).
Four operational details not stated in the plan (significance level $0.05$; ``holds at
$g \geq 4$'' as $p \leq 0.05$ at both $g = 4$ and $g = 8$; the decision based on the
largest $\Delta_{\min}$ at $g \geq 4$; averaging over subjects only for complete cells)
were committed before the primary run. An implementation error initially derived
$\Delta_{\min}$ from the primary analysis's own seed variability instead of the pilot's;
we report the pre-specified value and the original value as supplementary. The verdict
is the same under both.

\paragraph{Blinding and deviations.}
The primary analysis was not blind. Results for subject~1 at $g \in \{1, 4, 8\}$ and for
subjects~2 and~3 at $g = 1$ were known from earlier steps; only $g = 2$ and subjects~2
and~3 at $g \geq 4$ were new. Several pre-specified instruments were not implemented:
a capacity sweep, a PLS readout with 25 components, a ridge fit on the repeated-image
test set, and a leave-one-subject-out bound for the ridge scores. The image-wise cell
shuffle was run for one seed at 224\,px only, and the grid ladder only at 224\,px.

\paragraph{Exploratory analyses.}
All analyses in Sections~\ref{sec:standardization}--\ref{sec:gaingeometry} lie outside
the pre-specified plan, which excluded, among others, other layers and band-pass
decompositions. They were developed sequentially. For each step, predictions and
decision rules were written down and committed to version control before the step's
scripts were run; the commit timestamps, not the dates in file names, are authoritative
and are listed in the repository. We report every pre-stated prediction, including the
\FailedPredictions{} that were not met (Appendix~\ref{app:failed}). Run-grouped cross-validation
folds were an additional sensitivity analysis and were not part of the plan.

\paragraph{Inference and its limits.}
For RSA, we resampled stimuli (1000 bootstrap resamples, seed 20261002) and computed
differences between conditions from the same resamples; intervals are 95\% percentile
intervals. For ridge scores, seeds were the unit of inference. Intervals across the
exploratory analyses are not corrected for multiple comparisons. With three subjects,
all statements concern variation over stimuli and seeds and do not generalise to a
population of observers. Code, configuration, amendments and the full commit history
are available in the accompanying repository.
% TODO: repository URL after release
